\section{模型解释的数学形式化}
\label{sec:14-Deep-Learning/07-Interpretability-Mathematics/04}

你训练好了一个信用评分模型。输入是 30 个特征——收入、负债比、近期查询次数、账户年龄等等——输出是一个 0 到 1 之间的分数。一位客户被拒贷，你被要求回答：``哪个特征是这次决策的主要原因？''

你打开模型，有三种本能的做法。

第一种：看权重。线性模型的系数自然暗示了每个特征的重要程度——但你的模型是深度网络，根本没有每个特征的独立权重。

第二种：看梯度。对输入求导，梯度绝对值最大的特征似乎就是输出最敏感的。客户本次申请的``近期查询次数''梯度最大，你就报告查询次数是主要原因。

第三种：删掉特征重跑预测。把``近期查询次数''换成人群均值后，分数从 0.32 跳到了 0.61——变化很大，果然是它。

同一个客户，同一个模型，三种方法给出了三组不同的特征排名。梯度法把``账户年龄''排在第二，删除法把``收入''排在第二，而一位同事用 SHAP 跑出来的第一名是``负债比''。哪个是对的？

问题并不在于哪个方法计算错误。问题是``哪个特征重要''还没有被定义成一个数学问题。``重要''是相对于哪个输出？相对于哪个基线？通过哪类扰动来测量？在哪个数据分布上平均？在回答之前，你不知道自己问的究竟是什么。

\subsection{SHAP：每个特征是一个博弈参与者}

SHAP 把特征视为合作博弈中的参与者，目标是公平地分配模型输出与基线输出之间的差距——就像把合作产生的利润分摊给各个股东。它的数学核心是 Shapley 值。

给定一个集合函数 \(v(S)\)——它接受一个特征子集 \(S\)，返回仅使用这些特征时模型的输出（某种意义上的）——特征 \(i\) 的 Shapley 值定义为其边际贡献在所有可能加入顺序下的加权平均：

\[
\phi_i
= \sum_{S \subseteq N \setminus \{i\}}
\frac{|S|!\,(n - |S| - 1)!}{n!}
\bigl[v(S \cup \{i\}) - v(S)\bigr].
\]

权重 \(\frac{|S|!\,(n-|S|-1)!}{n!}\) 来自组合数学：在所有 \(n!\) 种逐个引入特征的顺序中，\(S\) 恰好是 \(i\) 之前的那个集合的顺序有多少种。Shapley 值是唯一满足四条公平性公理（效率、对称、虚拟参与者、可加性）的分配方案。

真正困难不在这个公式本身——它只是一个加权平均。困难在于定义 \(v(S)\)：模型需要全部特征才能运行，只给它一个子集 \(S\) 时，缺失的特征用什么值填充？

三种常见的填充策略回答了三种不同的问题。用边缘分布采样缺失特征——``在全体人群中，用特征 \(S\) 的已知值去推断输出，缺失特征按总体平均来看''。用条件分布采样——``在特征 \(S\) 取这些具体值的人群中，缺失特征通常取什么值''。用单个背景样本填充——``以这个特定参照人为基准，特征 \(S\) 的改变会导致多大差异''。

在特征独立时，前两种策略等价，且对线性模型化简为 \(w_i(x_i - \E[X_i])\)——系数乘以特征值偏离均值的量。但当特征相关时——比如收入和负债比负相关——不同填充策略会把相关特征之间的共享信息分配给不同的特征。SHAP 的归因结果因此不是模型本身的属性，而是``模型 + 特征分布 + 缺失值处理策略''这个三人博弈的属性。

\subsection{LIME：在局部训练一个替身来解释}

LIME 的思路完全不同。它不关心全局的公平分配公理，而是在某个具体样本的邻域内，用一个简单的、人类可读的模型（通常是稀疏线性模型）局部逼近原始模型的行为。

操作过程是：在目标样本附近按预设的扰动机制生成一批变体——比如随机遮蔽若干词、用背景值替换若干像素——用原模型给这些变体打分，再用距离加权的线性回归拟合这些分数。回归系数的绝对值就报告为特征重要性。

LIME 解释的是：在选定的邻域形状、扰动分布、核宽和替身模型类下，哪个简单函数最接近原模型在这个区域的行为。它不声称这是原模型的全局属性，也不声称系数反映了真实数据生成机制。

这里有三个隐蔽的自由度。邻域多大？太小了梯度近似刚好是线性，太大了线性替身可能根本拟合不了。扰动怎么生成？对图像做矩形遮挡生成的样本可能完全不像任何真实图像——这时替身主要在解释人造噪声上的行为，而不是数据流形上的行为。替身模型选什么类？选线性模型意味着只能捕获一阶依赖，选决策树能捕获交互但失去了系数排序的直觉。

如果模型足够光滑、邻域足够小、采样覆盖合理，线性替身的系数会趋近于局部梯度——此时 LIME 和基于梯度的归因方法收敛到同一个答案。一般情况下两者不必一致，报告 LIME 结果时也需要同时报告邻域定义、拟合误差和多次运行的稳定性。

\subsection{积分梯度：沿路径分配输出差异}

积分梯度从一个不同的完备性要求出发：输入 \(x\) 和某个基线 \(x_0\) 之间的输出差异，应该被全部分配给各个特征，不留残差。通过直线路径连接基线与输入，沿途累积每个特征的梯度贡献：

\[
\operatorname{IG}_i(x)
= (x_i - x_{0,i})
\int_0^1
\frac{\partial f(x_0 + \alpha(x - x_0))}{\partial x_i}
\,d\alpha.
\]

将全部坐标的 IG 求和：

\[
\sum_{i=1}^d \operatorname{IG}_i(x)
= \int_0^1 \sum_{i=1}^d
\frac{\partial f}{\partial x_i}
\cdot \frac{\partial (x_{0,i} + \alpha(x_i - x_{0,i}))}{\partial \alpha}
\,d\alpha
= f(x) - f(x_0).
\]

完备性（completeness）是这个定义的内置属性——不需要额外公理来保证。但它没有说明基线是否合理。基线表示``没有信息时的输出''，选全黑图像、全白图像、噪声图像或训练集均值，归因结果差异可能相当大。而直线路径是否落在数据支持集附近，积分梯度也同样不保证——路径可能穿越完全没有真实数据的区域，梯度在那些位置的行为对解释产生了影响，却说不出是从哪里来的。

饱和区是一个典型陷阱。当某个输入维度已经使网络中的某个神经元完全饱和时，该点上的梯度接近零，直接梯度法会报告该特征不重要。积分梯度通过累积路径上各点的梯度，能恢复出特征在到达饱和区之前做出的贡献。但恢复出来的贡献仍然描述的是模型在这个输入上的行为，不自动成为因果贡献：改变这个特征，模型输出会不会相应变化，积分梯度不回答这个问题。

\subsection{四种性质，解释方法至少需要接受其中几项的检验}

不同的解释即使各自内部一致，也需要有独立于其定义的评估手段：

\begin{itemize}
\tightlist
\item
  忠实度（fidelity）：解释能否预测模型在受控扰动下的输出变化？如果解释声称特征 A 最重要，那么只改变特征 A 时，输出的变化是否确实与解释预测的一致？
\item
  稳定性（stability）：对输入做微小扰动或更换随机种子后，解释是否无理由地剧烈变化？如果同一张图片连跑三次 SHAP 得到三组不同的前三名特征，解释就无法被信任为模型的稳定描述。
\item
  敏感性（sensitivity）：删除一个模型确实依赖的特征后，解释是否给出了相应的响应？一个永远报告``所有特征等权重要''的解释在所有输入上都稳定，但它不携带任何信息。
\item
  适用范围（scope）：结论是针对单个样本、局部邻域还是总体分布？SHAP 可以在全部样本上平均得到全局重要性，LIME 只在目标点邻域内有效，把一个样本的 LIME 解释当作全局结论是范围的误用。
\end{itemize}

删除高归因特征并观察性能下降，是最常用的评估手段。但这个做法有一个容易被忽略的陷阱：删除特征会制造分布外（out-of-distribution）输入。一个被训练来接收 30 个特征的模型，突然收到一个缺失了``收入''特征的样本——在它的训练经验里从未见过这种残缺输入——输出变化可能部分来自特征缺失的实际影响，部分来自模型在陌生输入上的任意行为。更好的评估会让扰动尽量满足数据生成约束，并设置已知真实机制的合成数据作为对照。

\subsection{归因、因果与机制是三个不同的问题}

SHAP、LIME 和积分梯度主要描述给定模型的输入依赖关系。它们回答的是``这个模型在做预测时用到了哪些输入信息''——归因（attribution）。

它们不直接回答：``改变现实世界中的特征 A，输出会怎样变化？''——这是因果（causation）问题，需要干预假设和识别策略，\hyperref[sec:13-Machine-Learning/12-Causal-Inference/02]{``结构因果模型把干预写成机制替换''}中讨论。

它们也不回答：``模型内部是否使用了人类赋予热力图的语义？''——这是机制（mechanism）问题，需要检验内部激活模式、回路行为或反事实干预在表示空间中的效果。

三个问题各有一组方法。把一个归因方法的热力图当作因果解释或机制证据，是范畴错误——即使热力图恰好与人类直觉吻合。吻合也许说明模型学到了合理的输入依赖，也许说明任务本身太简单以至于很多不同的依赖方式都能产生相似的预测。

选择解释方法之前，先写下四件事：被解释的是哪个输出值；参考对象或基线是什么；允许对输入做哪类扰动；结论将被谁用于做什么决定。数学形式化的价值不是让某一组数字看起来客观，而是让不同答案之间的差异来源变得可见——让解释可以被检验、被比较、被推翻。
